% page 198 of the facsimile (image p-197 of the scan)
% block 157.5 x 246.3 mm ; left margin 8.0 mm
\pgc{-12.700mm}{0.000mm}{
\l{\cel{3}190}
\l{}
\l{\sou{gazelo.(zool.}) Speco di antilopo qua habitas Afrika ed Azia.}
\l{\cel{3}L.\cel{1}\sou{antilope dorca}. - DEFIRS.}
\l{}
\l{\sou{gazolino}. (\sou{kemio)}. Petrol-etero, la elemento maxim volatila}
\l{\cel{3}di la petrolo kruda.}
\l{}
\l{\sou{gazono}. (\sou{bot.)}\cel{2}Herbo tre tenua ek graminei diversa, qua,}
\l{\cel{3}sur sulo, formacas quaza tapiso de verdajo. - F.}
\l{}
\l{\sou{ge-}.\cel{2}Prefixo qua signifikas \textquotedbl{}maskulo e femino, samasorta}
\l{\cel{3}egardata pare\textquotedbl{}.}
\l{}
\l{\sou{geblo}. (\sou{arkitekt.}) Parto supra di la muro, qua finas per pinto}
\l{\cel{3}qua portas la extremajo di la kolmo-trabo di la tekto-kar-}
\l{\cel{3}penturo. - DE.}
\l{}
\l{\sou{geheno}. (\sou{biblo}).Inferno. - DEFIRS.}
\l{}
\l{\sou{\textquotedbl{}geisha}\textquotedbl{}. Dansistino e kantistino Japoniana.}
\l{}
\l{\sou{gelatino}. Substanco azotoza quan onu extraktas de la osti,de}
\l{\cel{3}la kartilagi di la animali, e qua formacas, kun aquo, je-}
\l{\cel{3}teo. - DEFIRS.}
\l{}
\l{\sou{gelto}. Parto proporcionala de la pekunio qua rezultas de la}
\l{\cel{3}vendo di la vari, grantata a sua komizi (ultre la salario-}
\l{\cel{3}quanto fixa di ici) da ula butikegi. - DF.}
\l{}
\l{\sou{gemo}. Lapido tre precoza. - DEFIS.}
\l{}
\l{\sou{gemulo}. (\sou{bot.)}\cel{2}I. Rudimento di la stipo, qua aparas lor la}
\l{\cel{3}jermifo. - II. Mikra korpo ovoida qua, separante su de la}
\l{\cel{3}planto, destinesas a reproduktar ica.}
\l{}
\l{\sou{genciano}. (\sou{bot.})\cel{2}Planto de la familio \textquotedbl{}genciani\textquotedbl{}, di qua la}
\l{\cel{3}suko esas bitra, qua kreskas en la monti. - L.\cel{1}\sou{gentiana}.}
\l{\cel{3}- DEFIRS.}
\l{}
\l{\sou{genealogio}. Decendo-lineo pri un o plura personi, quan pro-}
\l{\cel{3}duktas la suceso da lua o lia ancestri. - DEFIRS.}
\l{}
\l{\sou{genero}. I. Grupo naturala de enti qui intersimilesas pri ula}
\l{\cel{3}traiti esencala. - II. Ideo generala sub qua la spirito}
\l{\cel{3}asemblas to quo esas komuna ye termini diversa, per elimi-}
\l{\cel{3}nar la difero-punti. - III. Ensemblo de la karakteri esen-}
\l{\cel{3}cala di kozo. - DEFIS.}
\l{}
\l{\sou{generaciono}. La homi qui vivas en la sama periodo (evaluata}
\l{\cel{3}segun la meza duro di la homo-vivo). - DEFIRS.}
\l{}
\l{\sou{generalo}. (\sou{armeo)}. Milito-chefo qua komandas armeo, diviziono}
\l{\cel{3}o brigado. - DEFIRS.}
\l{}
\l{\sou{generalisimo.}\cel{2}(\sou{armeo)}. La generalo qua esas la chefo di omna}
\l{\cel{3}generali cetera. - DEFIS.}
}
